\begin{PSbox}[unbreakable]{obj}{Learning Objectives}

After completing this module, students will be able to:

\begin{enumerate}
\def\labelenumi{\arabic{enumi}.}
\tightlist
\item
  Compute and interpret covariance
\item
  Compute and interpret the correlation coefficient
\item
  Distinguish between independence and uncorrelatedness
\item
  Construct and interpret correlation matrices
\item
  Apply correlation analysis to engineering problems
\end{enumerate}

\end{PSbox}

\PSsection{9.1}{Introduction: Measuring Linear Relationships}

When two random variables are not independent, we need to quantify \textbf{how they relate}. Covariance and correlation measure the \textbf{linear relationship} between two RVs.

Engineering examples:

\begin{itemize}
\tightlist
\item
  Correlated noise in adjacent sensor channels
\item
  Signal measurements at two antennas (MIMO correlation)
\item
  Voltage and current in a circuit (related through impedance)
\item
  Temperature and device performance (thermal dependence)
\end{itemize}

\PSsection{9.2}{Covariance}

\begin{PSbox}[unbreakable]{def}{Definition}

\PSdisplay{\text{Cov}(X,Y) = E[(X - \mu_X)(Y - \mu_Y)] = E[XY] - E[X]E[Y]}

\end{PSbox}

\PSsubsection{Properties}

\begin{enumerate}
\def\labelenumi{\arabic{enumi}.}
\tightlist
\item
  $\operatorname{Cov}(X,\, X) = \operatorname{Var}(X)$
\item
  $\operatorname{Cov}(X,\, Y) = \operatorname{Cov}(Y,\, X)$ (symmetric)
\item
  $\operatorname{Cov}(\text{aX} + b,\, \text{cY} + d)$ = ac $\cdot \operatorname{Cov}(X,\, Y)$
\item
  If $X,\, Y$ independent $\to \operatorname{Cov}(X,\, Y) = 0$
\item
  $\operatorname{Var}(X + Y) = \operatorname{Var}(X) + \operatorname{Var}(Y) + 2\operatorname{Cov}(X,\, Y)$
\item
  $\operatorname{Var}(X - Y) = \operatorname{Var}(X) + \operatorname{Var}(Y) - 2\operatorname{Cov}(X,\, Y)$
\end{enumerate}

\PSsubsection{Sign Interpretation}

{\def\LTcaptype{none} % do not increment counter
\begin{longtable}[c]{@{}cc@{}}
\toprule\noalign{}
\rowcolor{PSnavy}\PSth{$\operatorname{Cov}(X,\,Y)$} & \PSth{Meaning} \\
\midrule\noalign{}
\endhead
\bottomrule\noalign{}
\endlastfoot
\textgreater{} 0 & $X$ and $Y$ tend to increase together \\
\textless{} 0 & When $X$ increases, $Y$ tends to decrease \\
= 0 & No linear relationship (uncorrelated) \\
\end{longtable}
}

\begin{PSbox}[unbreakable]{ex}{Engineering Example: Amplifier Gain and Output}

If gain $G$ varies randomly and input is fixed at $V_{\mathrm{in}}$:

\begin{itemize}
\tightlist
\item
  Output: $V_{\mathrm{out}} = G \cdot V_{\mathrm{in}}$
\item
  $\operatorname{Cov}(G,\, V_{\text{out}}) = \operatorname{Cov}(G,\, G \cdot V_{\mathrm{in}}) = V_{\mathrm{in}} \cdot \operatorname{Var}(G) > 0$
\end{itemize}

Higher gain \ensuremath{\to} higher output (positive covariance).

\end{PSbox}

\PSsection{9.3}{Correlation Coefficient}

\begin{PSbox}[unbreakable]{def}{Definition}

\PSdisplay{\rho_{XY} = \frac{\text{Cov}(X,Y)}{\sigma_X \sigma_Y}}

\end{PSbox}

\PSsubsection{Properties}

\begin{enumerate}
\def\labelenumi{\arabic{enumi}.}
\tightlist
\item
  \textbf{Bounded:} $-1 \le \rho \le 1$
\item
  \textbf{$\rho = 1$:} Perfect positive linear relationship ($Y$ = aX $+ b,\, a > 0$)
\item
  \textbf{$\rho = -1$:} Perfect negative linear relationship ($Y$ = aX $+ b,\, a < 0$)
\item
  \textbf{$\rho = 0$:} Uncorrelated (no linear relationship)
\item
  \textbf{Dimensionless:} Unlike covariance, $\rho$ doesn’t depend on units
\end{enumerate}

\PSsubsection{Interpretation Scale}

{\def\LTcaptype{none} % do not increment counter
\begin{longtable}[c]{@{}cc@{}}
\toprule\noalign{}
\rowcolor{PSnavy}& \PSth{$\rho$} \\
\midrule\noalign{}
\endhead
\bottomrule\noalign{}
\endlastfoot
0.0 – 0.2 & Very weak/negligible \\
0.2 – 0.4 & Weak \\
0.4 – 0.6 & Moderate \\
0.6 – 0.8 & Strong \\
0.8 – 1.0 & Very strong \\
\end{longtable}
}

\begin{PSbox}[unbreakable]{ex}{Engineering Example: Correlated Fading}

Two antenna paths with correlation $\rho = 0.3$:

\begin{itemize}
\tightlist
\item
  Low correlation \ensuremath{\to} diversity works well (independent fading)
\item
  High correlation \ensuremath{\to} diversity gain is reduced
\end{itemize}

MIMO system design requires: $\rho < 0.5$ for effective spatial multiplexing.

\end{PSbox}

\PSsection{9.4}{Cross-Correlation}

\begin{PSbox}[unbreakable]{def}{Definition}

\PSdisplay{R_{XY} = E[XY]}

\end{PSbox}

\PSsubsection{Relationship to Covariance}

\PSdisplay{\text{Cov}(X,Y) = R_{XY} - \mu_X \mu_Y}

If either $X$ or $Y$ has zero mean: $\operatorname{Cov}(X,\,Y) = E[XY] = R_{XY}$

\PSsubsection{Engineering Significance}

Cross-correlation is fundamental in:

\begin{itemize}
\tightlist
\item
  Matched filters (correlation receivers in communications)
\item
  Radar signal processing (detecting time delay)
\item
  Signal detection (correlating received with known template)
\end{itemize}

\PSsection{9.5}{Independence vs. Uncorrelatedness}

\PSsubsection{The Critical Distinction}

{\def\LTcaptype{none} % do not increment counter
\begin{longtable}[c]{@{}cc@{}}
\toprule\noalign{}
\rowcolor{PSnavy}\PSth{Statement} & \PSth{Always True?} \\
\midrule\noalign{}
\endhead
\bottomrule\noalign{}
\endlastfoot
Independent \ensuremath{\to} Uncorrelated & \PSyes{} YES \\
Uncorrelated \ensuremath{\to} Independent & \PSno{} NO (in general) \\
Uncorrelated \ensuremath{\to} Independent (jointly Gaussian) & \PSyes{} YES \\
\end{longtable}
}

\PSsubsection{Why Independent \ensuremath{\to} Uncorrelated}

If $X,\, Y$ independent: $E[XY] = E[X]E[Y]$, so $\operatorname{Cov}(X,\,Y) = E[XY] - E[X]E[Y] = 0$.

\PSsubsection{Counterexample: Uncorrelated but Dependent}

Let $X \sim N(0,\,1)$ and $Y = X^{2}$. Then:

\begin{itemize}
\tightlist
\item
  $\operatorname{Cov}(X,\, Y) = E[XY] - E[X]E[Y] = E[X^{3}] - 0 = 0$ (since $X$ is symmetric)
\item
  But $Y$ is completely determined by $X$ \ensuremath{\to} maximally dependent!
\end{itemize}

\textbf{$\rho = 0$ does NOT mean independent} (except for jointly Gaussian RVs).

\PSsubsection{The Gaussian Exception}

For \textbf{jointly Gaussian} random variables:

\PSdisplay{\text{Uncorrelated} \Leftrightarrow \text{Independent}}

This is why the Gaussian distribution is so convenient in engineering — checking $\rho = 0$ suffices to establish independence.

\PSsubsection{Engineering Implication}

In communications with Gaussian noise:

\begin{itemize}
\tightlist
\item
  If noise samples are uncorrelated \ensuremath{\to} they are independent
\item
  This simplifies analysis enormously
\item
  For non-Gaussian interference: uncorrelated does NOT guarantee independence
\end{itemize}

\PSsection{9.6}{Correlation Matrix}

\begin{PSbox}[unbreakable]{def}{Definition for Random Vector $X = [X_{1},\, X_{2},\, \ldots,\, X_{n}]^{T}$}

The \textbf{covariance matrix} $\Sigma$:

\PSdisplay{\Sigma_{ij} = \text{Cov}(X_i, X_j)}

\PSdisplay{\Sigma = \begin{bmatrix} \text{Var}(X_1) & \text{Cov}(X_1,X_2) & \cdots \\ \text{Cov}(X_2,X_1) & \text{Var}(X_2) & \cdots \\ \vdots & & \ddots \end{bmatrix}}

The \textbf{correlation matrix} $R$ (normalized):

\PSdisplay{R_{ij} = \rho_{X_i X_j} = \frac{\Sigma_{ij}}{\sqrt{\Sigma_{ii}\Sigma_{jj}}}}

\end{PSbox}

\PSsubsection{Properties of Covariance Matrix}

\begin{enumerate}
\def\labelenumi{\arabic{enumi}.}
\tightlist
\item
  Symmetric: $\Sigma = \Sigma^{T}$
\item
  Positive semi-definite: $x^{T}\Sigma x \ge 0$ for all $x$
\item
  Diagonal entries = variances
\item
  Off-diagonal entries = covariances
\end{enumerate}

\begin{PSbox}[unbreakable]{ex}{Engineering Example: Three Sensors}

Three temperature sensors with:

\begin{itemize}
\tightlist
\item
  $\sigma_{1} = \sigma_{2} = \sigma_{3} = 1^{\circ}C$
\item
  $\rho_{12} = 0.8$ (sensors 1,2 are nearby)
\item
  $\rho_{13} = 0.3$ (sensor 3 is farther)
\item
  $\rho_{23} = 0.4$
\end{itemize}

Correlation matrix:\PSdisplay{R = \begin{bmatrix} 1.0 & 0.8 & 0.3 \\ 0.8 & 1.0 & 0.4 \\ 0.3 & 0.4 & 1.0 \end{bmatrix}}

\end{PSbox}

\PSsection{9.7}{Variance of Linear Combinations}

\PSsubsection{General Formula}

\PSdisplay{\text{Var}\left(\sum_{i=1}^n a_i X_i\right) = \sum_{i=1}^n a_i^2 \text{Var}(X_i) + 2\sum_{i<j} a_i a_j \text{Cov}(X_i, X_j)}

In matrix form: $\operatorname{Var}(a^{T}X) = a^{T}\Sigma a$

\PSsubsection{Special Case: Sum of Two RVs}

$\operatorname{Var}(X + Y) = \operatorname{Var}(X) + \operatorname{Var}(Y) + 2\operatorname{Cov}(X,\,Y) = \sigma_{X}^{2} + \sigma_{Y}^{2} + 2\rho \sigma_{X}\sigma_{Y}$

{\def\LTcaptype{none} % do not increment counter
\begin{longtable}[c]{@{}
  >{\centering\arraybackslash}p{(0.965\linewidth - 4\tabcolsep) * \real{0.4194}}
  >{\centering\arraybackslash}p{(0.965\linewidth - 4\tabcolsep) * \real{0.3226}}
  >{\centering\arraybackslash}p{(0.965\linewidth - 4\tabcolsep) * \real{0.2581}}@{}}
\toprule\noalign{}
\rowcolor{PSnavy}\PSth{Correlation} & \PSth{$\operatorname{Var}(X+Y)$} & \PSth{Effect} \\
\midrule\noalign{}
\endhead
\bottomrule\noalign{}
\endlastfoot
$\rho = 1$ & $(\sigma_{X} + \sigma_{Y})^{2}$ & Maximum — variances add constructively \\
$\rho = 0$ & $\sigma_{X}^{2} + \sigma_{Y}^{2}$ & Independent addition \\
$\rho = -1$ & $(\sigma_{X} - \sigma_{Y})^{2}$ & Minimum — cancellation \\
\end{longtable}
}

\PSsubsection{Engineering Application: Averaging Correlated Sensors}

If $n$ sensors have equal variance $\sigma^{2}$ and pairwise correlation $\rho$:

\PSdisplay{\text{Var}(\bar{X}) = \frac{\sigma^2}{n}[1 + (n-1)\rho]}

\begin{itemize}
\tightlist
\item
  If $\rho = 0$ (independent): $\operatorname{Var}(\bar{X}) = \frac{\sigma^{2}}{n}$ (full averaging benefit)
\item
  If $\rho = 1$ (identical): $\operatorname{Var}(\bar{X}) = \sigma^{2}$ (no benefit from averaging!)
\item
  If $\rho = 0.5,\, n = 4$: $\operatorname{Var}(\bar{X}) = \frac{\sigma^{2}(1 + 3 \times 0.5)}{4} = 0.625\sigma^{2}$ (reduced benefit)
\end{itemize}

\PSsection{9.8}{MATLAB Examples}

\begin{PSbox}[unbreakable]{ex}{Example 1: Computing Correlation}

\tcbset{PScodemode/.style={unbreakable}}

\begin{Shaded}
\begin{Highlighting}[]
\CommentTok{\%\% Sample Correlation: Correlated Sensor Data}
\VariableTok{N} \OperatorTok{=} \FloatTok{10000}\OperatorTok{;}
\VariableTok{rho} \OperatorTok{=} \FloatTok{0.7}\OperatorTok{;}  \VariableTok{sigma1} \OperatorTok{=} \FloatTok{2}\OperatorTok{;}  \VariableTok{sigma2} \OperatorTok{=} \FloatTok{3}\OperatorTok{;}

\CommentTok{\% Generate correlated Gaussians}
\VariableTok{Z1} \OperatorTok{=} \VariableTok{randn}\NormalTok{(}\FloatTok{1}\OperatorTok{,} \VariableTok{N}\NormalTok{)}\OperatorTok{;}
\VariableTok{Z2} \OperatorTok{=} \VariableTok{randn}\NormalTok{(}\FloatTok{1}\OperatorTok{,} \VariableTok{N}\NormalTok{)}\OperatorTok{;}
\VariableTok{X} \OperatorTok{=} \VariableTok{sigma1} \OperatorTok{*} \VariableTok{Z1}\OperatorTok{;}
\VariableTok{Y} \OperatorTok{=} \VariableTok{sigma2} \OperatorTok{*}\NormalTok{ (}\VariableTok{rho}\OperatorTok{*}\VariableTok{Z1} \OperatorTok{+} \VariableTok{sqrt}\NormalTok{(}\FloatTok{1}\OperatorTok{{-}}\VariableTok{rho}\OperatorTok{\^{}}\FloatTok{2}\NormalTok{)}\OperatorTok{*}\VariableTok{Z2}\NormalTok{)}\OperatorTok{;}

\CommentTok{\% Compute sample statistics}
\VariableTok{cov\_xy} \OperatorTok{=} \VariableTok{mean}\NormalTok{(}\VariableTok{X}\OperatorTok{.*}\VariableTok{Y}\NormalTok{) }\OperatorTok{{-}} \VariableTok{mean}\NormalTok{(}\VariableTok{X}\NormalTok{)}\OperatorTok{*}\VariableTok{mean}\NormalTok{(}\VariableTok{Y}\NormalTok{)}\OperatorTok{;}
\VariableTok{rho\_hat} \OperatorTok{=} \VariableTok{cov\_xy} \OperatorTok{/}\NormalTok{ (}\VariableTok{std}\NormalTok{(}\VariableTok{X}\NormalTok{)}\OperatorTok{*}\VariableTok{std}\NormalTok{(}\VariableTok{Y}\NormalTok{))}\OperatorTok{;}

\VariableTok{fprintf}\NormalTok{(}\SpecialStringTok{\textquotesingle{}True ρ = \%.2f, Estimated ρ = \%.4f\textbackslash{}n\textquotesingle{}}\OperatorTok{,} \VariableTok{rho}\OperatorTok{,} \VariableTok{rho\_hat}\NormalTok{)}\OperatorTok{;}
\VariableTok{fprintf}\NormalTok{(}\SpecialStringTok{\textquotesingle{}Cov(X,Y): Theory=\%.2f, Estimated=\%.4f\textbackslash{}n\textquotesingle{}}\OperatorTok{,} \VariableTok{rho}\OperatorTok{*}\VariableTok{sigma1}\OperatorTok{*}\VariableTok{sigma2}\OperatorTok{,} \VariableTok{cov\_xy}\NormalTok{)}\OperatorTok{;}

\CommentTok{\% Scatter plot}
\VariableTok{figure}\OperatorTok{;}
\VariableTok{scatter}\NormalTok{(}\VariableTok{X}\OperatorTok{,} \VariableTok{Y}\OperatorTok{,} \FloatTok{1}\OperatorTok{,} \SpecialStringTok{\textquotesingle{}b\textquotesingle{}}\OperatorTok{,} \SpecialStringTok{\textquotesingle{}.\textquotesingle{}}\NormalTok{)}\OperatorTok{;}
\VariableTok{xlabel}\NormalTok{(}\SpecialStringTok{\textquotesingle{}X (Sensor 1)\textquotesingle{}}\NormalTok{)}\OperatorTok{;} \VariableTok{ylabel}\NormalTok{(}\SpecialStringTok{\textquotesingle{}Y (Sensor 2)\textquotesingle{}}\NormalTok{)}\OperatorTok{;}
\VariableTok{title}\NormalTok{(}\VariableTok{sprintf}\NormalTok{(}\SpecialStringTok{\textquotesingle{}Correlated Sensors (ρ = \%.2f)\textquotesingle{}}\OperatorTok{,} \VariableTok{rho}\NormalTok{))}\OperatorTok{;}
\VariableTok{grid} \VariableTok{on}\OperatorTok{;}
\end{Highlighting}
\end{Shaded}

\end{PSbox}

\begin{PSbox}[unbreakable]{ex}{Example 2: Uncorrelated but Dependent}

\tcbset{PScodemode/.style={unbreakable}}

\begin{Shaded}
\begin{Highlighting}[]
\CommentTok{\%\% Demonstration: ρ = 0 does NOT mean independent}
\VariableTok{N} \OperatorTok{=} \FloatTok{100000}\OperatorTok{;}
\VariableTok{X} \OperatorTok{=} \VariableTok{randn}\NormalTok{(}\FloatTok{1}\OperatorTok{,} \VariableTok{N}\NormalTok{)}\OperatorTok{;}
\VariableTok{Y} \OperatorTok{=} \VariableTok{X}\OperatorTok{.\^{}}\FloatTok{2}\OperatorTok{;}        \CommentTok{\% Y completely depends on X}

\VariableTok{rho\_hat} \OperatorTok{=} \VariableTok{corr}\NormalTok{(}\VariableTok{X}\OperatorTok{\textquotesingle{},} \VariableTok{Y}\OperatorTok{\textquotesingle{}}\NormalTok{)}\OperatorTok{;}
\VariableTok{fprintf}\NormalTok{(}\SpecialStringTok{\textquotesingle{}Correlation ρ(X, X²) = \%.4f ≈ 0 (UNCORRELATED)\textbackslash{}n\textquotesingle{}}\OperatorTok{,} \VariableTok{rho\_hat}\NormalTok{)}\OperatorTok{;}
\VariableTok{fprintf}\NormalTok{(}\SpecialStringTok{\textquotesingle{}But Y = X² is perfectly determined by X (DEPENDENT!)\textbackslash{}n\textquotesingle{}}\NormalTok{)}\OperatorTok{;}

\VariableTok{figure}\OperatorTok{;}
\VariableTok{scatter}\NormalTok{(}\VariableTok{X}\OperatorTok{,} \VariableTok{Y}\OperatorTok{,} \FloatTok{1}\OperatorTok{,} \SpecialStringTok{\textquotesingle{}.\textquotesingle{}}\NormalTok{)}\OperatorTok{;} \VariableTok{xlabel}\NormalTok{(}\SpecialStringTok{\textquotesingle{}X\textquotesingle{}}\NormalTok{)}\OperatorTok{;} \VariableTok{ylabel}\NormalTok{(}\SpecialStringTok{\textquotesingle{}Y = X²\textquotesingle{}}\NormalTok{)}\OperatorTok{;}
\VariableTok{title}\NormalTok{(}\SpecialStringTok{\textquotesingle{}Uncorrelated but Dependent!\textquotesingle{}}\NormalTok{)}\OperatorTok{;}
\end{Highlighting}
\end{Shaded}

\end{PSbox}

\begin{PSbox}[unbreakable]{ex}{Example 3: Correlation Matrix Visualization}

\tcbset{PScodemode/.style={unbreakable}}

\begin{Shaded}
\begin{Highlighting}[]
\CommentTok{\%\% Correlation Matrix for Multiple Sensors}
\VariableTok{n\_sensors} \OperatorTok{=} \FloatTok{5}\OperatorTok{;}
\CommentTok{\% Create correlation matrix with decaying correlation}
\VariableTok{R} \OperatorTok{=} \VariableTok{zeros}\NormalTok{(}\VariableTok{n\_sensors}\NormalTok{)}\OperatorTok{;}
\KeywordTok{for} \VariableTok{i} \OperatorTok{=} \FloatTok{1}\OperatorTok{:}\VariableTok{n\_sensors}
    \KeywordTok{for} \VariableTok{j} \OperatorTok{=} \FloatTok{1}\OperatorTok{:}\VariableTok{n\_sensors}
        \VariableTok{R}\NormalTok{(}\VariableTok{i}\OperatorTok{,}\VariableTok{j}\NormalTok{) }\OperatorTok{=} \FloatTok{0.9}\OperatorTok{\^{}}\VariableTok{abs}\NormalTok{(}\VariableTok{i}\OperatorTok{{-}}\VariableTok{j}\NormalTok{)}\OperatorTok{;}  \CommentTok{\% Correlation decays with distance}
    \KeywordTok{end}
\KeywordTok{end}

\CommentTok{\% Generate correlated samples}
\VariableTok{L} \OperatorTok{=} \VariableTok{chol}\NormalTok{(}\VariableTok{R}\OperatorTok{,} \SpecialStringTok{\textquotesingle{}lower\textquotesingle{}}\NormalTok{)}\OperatorTok{;}   \CommentTok{\% Cholesky decomposition}
\VariableTok{N} \OperatorTok{=} \FloatTok{10000}\OperatorTok{;}
\VariableTok{Z} \OperatorTok{=} \VariableTok{randn}\NormalTok{(}\VariableTok{n\_sensors}\OperatorTok{,} \VariableTok{N}\NormalTok{)}\OperatorTok{;}
\VariableTok{X} \OperatorTok{=} \VariableTok{L} \OperatorTok{*} \VariableTok{Z}\OperatorTok{;}              \CommentTok{\% Correlated samples}

\CommentTok{\% Estimated correlation matrix}
\VariableTok{R\_hat} \OperatorTok{=} \VariableTok{corrcoef}\NormalTok{(}\VariableTok{X}\OperatorTok{\textquotesingle{}}\NormalTok{)}\OperatorTok{;}

\VariableTok{figure}\OperatorTok{;}
\VariableTok{subplot}\NormalTok{(}\FloatTok{1}\OperatorTok{,}\FloatTok{2}\OperatorTok{,}\FloatTok{1}\NormalTok{)}\OperatorTok{;} \VariableTok{imagesc}\NormalTok{(}\VariableTok{R}\NormalTok{)}\OperatorTok{;} \VariableTok{colorbar}\OperatorTok{;}
\VariableTok{title}\NormalTok{(}\SpecialStringTok{\textquotesingle{}Theoretical R\textquotesingle{}}\NormalTok{)}\OperatorTok{;} \VariableTok{xlabel}\NormalTok{(}\SpecialStringTok{\textquotesingle{}Sensor\textquotesingle{}}\NormalTok{)}\OperatorTok{;} \VariableTok{ylabel}\NormalTok{(}\SpecialStringTok{\textquotesingle{}Sensor\textquotesingle{}}\NormalTok{)}\OperatorTok{;}
\VariableTok{subplot}\NormalTok{(}\FloatTok{1}\OperatorTok{,}\FloatTok{2}\OperatorTok{,}\FloatTok{2}\NormalTok{)}\OperatorTok{;} \VariableTok{imagesc}\NormalTok{(}\VariableTok{R\_hat}\NormalTok{)}\OperatorTok{;} \VariableTok{colorbar}\OperatorTok{;}
\VariableTok{title}\NormalTok{(}\SpecialStringTok{\textquotesingle{}Estimated R\textquotesingle{}}\NormalTok{)}\OperatorTok{;} \VariableTok{xlabel}\NormalTok{(}\SpecialStringTok{\textquotesingle{}Sensor\textquotesingle{}}\NormalTok{)}\OperatorTok{;} \VariableTok{ylabel}\NormalTok{(}\SpecialStringTok{\textquotesingle{}Sensor\textquotesingle{}}\NormalTok{)}\OperatorTok{;}
\end{Highlighting}
\end{Shaded}

\end{PSbox}

\PSsection{9.9}{Practice Problems}

\begin{PSbox}[unbreakable]{prob}{Problem 1}

$X$ and $Y$ have: $E[X] = 2,\, E[Y] = 3,\, \operatorname{Var}(X) = 4,\, \operatorname{Var}(Y) = 9,\, E[XY] = 8$.

\begin{enumerate}
\def\labelenumi{(\alph{enumi})}
\tightlist
\item
  Find $\operatorname{Cov}(X,\,Y)$. (b) Find $\rho$. (c) Find $\operatorname{Var}(X+Y)$. (d) Find $\operatorname{Var}(2X-3Y)$.
\end{enumerate}

\PSsolution{} 

\begin{enumerate}
\def\labelenumi{(\alph{enumi})}
\tightlist
\item
  $\operatorname{Cov}(X,\,Y) = E[XY] - E[X]E[Y] = 8 - 2(3) = 2$
\item
  $\displaystyle \rho = \frac{2}{\sqrt{4} \cdot \sqrt{9}} = \frac{2}{6} = \frac{1}{3}$
\item
  $\operatorname{Var}(X+Y) = 4 + 9 + 2(2) = 17$
\item
  $\operatorname{Var}(2X-3Y) = 4(4) + 9(9) + 2(2)(-3)(2) = 16 + 81 - 24 = 73$
\end{enumerate}

\end{PSbox}

\begin{PSbox}[unbreakable]{prob}{Problem 2}

Four identical sensors $(\sigma^{2} = 1)$ have pairwise correlation $\rho = 0.5$. What is $\operatorname{Var}(\bar{X})$?

\PSsolution{} $\displaystyle \operatorname{Var}(\bar{X}) = \frac{\sigma^{2}}{n} \cdot [1 + (n-1)\rho] = \frac{1}{4} \cdot [1 + 3(0.5)] = 0.25 \times 2.5 = 0.625$

Compare: if independent $(\rho = 0)$: $\operatorname{Var}(\bar{X}) = 0.25$. Correlation reduces the effectiveness of averaging.

\end{PSbox}

\begin{PSbox}[unbreakable]{prob}{Problem 3}

Given $X \sim N(0,\,1)$, define $Y = X$ when $\lvert X\rvert < 1$ and $Y = -X$ when $\lvert X\rvert \ge 1$. Show that $\operatorname{Cov}(X,\,Y) \ne 0$ in this case but that $E[X \cdot Y]$ can be computed.

\PSsolution{} $E[XY] = E[X^{2} \cdot 1(\lvert X\rvert < 1)] + E[-X^{2} \cdot 1(\lvert X\rvert \ge 1)] = E[X^{2} \cdot 1(\lvert X\rvert < 1)] - E[X^{2} \cdot 1(\lvert X\rvert \ge 1)]$\PSbr{}$= P(\lvert X\rvert < 1)E[X^{2}\lvert \rvert X \mid < 1] - P(\lvert X\rvert \ge 1)E[X^{2}\lvert \rvert X \mid \ge 1]$\PSbr{}This is nonzero since $E[X^{2}\lvert \rvert X \mid < 1] \ne E[X^{2}\lvert \rvert X \mid \ge 1]$, so $\operatorname{Cov}(X,\,Y) = E[XY] - 0 \ne 0$.\PSbr{}$X$ and $Y$ are clearly dependent ($Y$ is defined from $X$), and in this case also correlated.

\emph{Next Module: {Module 10 — Multiple Random Variables}}

\end{PSbox}
